% Research fragment; no document preamble. All labels are prefixed gs:.
% Derived at repository commit fc4d3bf80534ad7c901d3b8c9e71baf2df0064ed.

\section{Gauss summation at zero balance and Stieltjes jets}
\label{gs:section}

The circular correlation identities in the incoming manuscripts concern
products of periodic Hurwitz functions.  The present section develops a
different, discrete construction: a gamma-ratio coefficient is compared
term by term with its leading Hurwitz-zeta coefficient.  At zero balance,
this subtraction converts the pole in Gauss summation into an ordinary,
absolutely convergent series.  Its complete parameter jet contains
generalized Stieltjes constants, polygamma functions, and nonlinear
generalized harmonic sums.

Gauss summation, gamma-ratio asymptotics, and the use of hypergeometric
parameter derivatives to evaluate harmonic sums are classical
\cite{gs:DLMF,gs:ChenChu,gs:Choi}.  Hypergeometric representations of
Stieltjes constants also have an established literature
\cite{gs:Coffey}.  The contribution here is an explicit common
normalization at zero balance, an every-index mixed-jet statement, its
polylogarithmic boundary interpretation, and concrete corrected formulas
for a group of historically printed conjectures.  No claim of first
discovery of Gauss differentiation, or of the present-day open status of
those conjectures, is made.

\subsection{A convergent master identity}

Use the convention
\begin{equation}\label{gs:stieltjes-convention}
 \zeta(1+s,c)=\frac1s+
 \sum_{k=0}^{\infty}\frac{(-1)^k\gamma_k(c)}{k!}s^k.
\end{equation}
Initially $a,b,c>0$.  All statements below extend holomorphically to
$\Re a,\Re b,\Re c>0$, with the principal logarithms of $n+c$.
Define
\begin{align}
 W_n(a,b;s)&=
 \frac{\Gamma(n+a)\Gamma(n+b)}
 {\Gamma(n+a+b+s)\Gamma(n+1)},\label{gs:weight}\\
 R_{a,b}(s)&=
 \frac{\Gamma(1+s)\Gamma(a)\Gamma(b)}
 {\Gamma(a+s)\Gamma(b+s)}.\label{gs:R}
\end{align}
Reciprocal gamma factors are interpreted as entire functions; a zero
of such a factor is not an excluded parameter value.

\begin{theorem}[Gauss--Stieltjes subtraction]\label{gs:master}
For $\Re a,\Re b,\Re c>0$ and $\Re s>-1$, the series
\begin{equation}\label{gs:series}
 \mathcal F_{a,b,c}(s)=
 \sum_{n=0}^{\infty}
 \left(W_n(a,b;s)-(n+c)^{-1-s}\right)
\end{equation}
converges normally on compact subsets of the full parameter domain and
is jointly holomorphic.  It satisfies
\begin{equation}\label{gs:master-equation}
 \mathcal F_{a,b,c}(s)
 =\frac{R_{a,b}(s)-1}{s}
  -\left(\zeta(1+s,c)-\frac1s\right).
\end{equation}
Both quotients on the right have their removable values at $s=0$.
Consequently, every finite mixed derivative in $a,b,c,s$ may be taken
term by term in \eqref{gs:series}.
\end{theorem}

\begin{proof}
The gamma-ratio expansion gives, uniformly on any compact subset of
the parameter domain,
\[
 W_n(a,b;s)=n^{-1-s}\bigl(1+O(n^{-1})\bigr),\qquad
 (n+c)^{-1-s}=n^{-1-s}\bigl(1+O(n^{-1})\bigr).
\]
For precision concerning uniformity, apply the logarithmic Stirling
expansion on a fixed sector containing all sufficiently large $n+a$,
$n+b$, $n+a+b+s$, and $n+1$.  On a compact parameter set the shifts are
bounded, and the remainder after its first correction is uniform.
Exponentiating gives the displayed estimates; this is the standard
gamma-ratio expansion recorded in DLMF~5.11.13.
If $\Re s\ge -1+\eta$ on the compact set, with $\eta>0$, the difference
is $O(n^{-1-\eta})$.  Thus the series converges normally and defines a
jointly holomorphic function.  Cauchy's formula on slightly larger
compact polydiscs justifies all finite parameter derivatives.

For $\Re s>0$, both series can be summed separately.  Gauss summation,
DLMF~15.4.20, gives
\[
 \sum_{n=0}^{\infty}W_n(a,b;s)
 =\frac{\Gamma(a)\Gamma(b)}{\Gamma(a+b+s)}
 {}_2F_1(a,b;a+b+s;1)
 =\frac{\Gamma(a)\Gamma(b)\Gamma(s)}
 {\Gamma(a+s)\Gamma(b+s)}.
\]
Subtract $\zeta(1+s,c)$ and use $\Gamma(1+s)=s\Gamma(s)$.
Since $R_{a,b}(0)=1$, both sides extend holomorphically through zero.
The identity theorem in $s$ proves \eqref{gs:master-equation} throughout
$\Re s>-1$ for every permitted $a,b,c$.  Joint holomorphy has already
been established by normal convergence.
\end{proof}

This is an ordinary convergent difference series.  Splitting the two
summands into separately divergent series at $s=0$ would lose the
specified normalization.  The arbitrary comparison shift $c$ is useful:
the gamma-ratio coefficient does not change, while its logarithmic
comparison term selects the generalized constant $\gamma_k(c)$.

\subsection{Every spectral derivative in Bell-polynomial form}

Let $Y_k$ denote the complete exponential Bell polynomial, defined by
\begin{equation}\label{gs:Bell-definition}
 \exp\left(\sum_{j\ge1}x_j\frac{t^j}{j!}\right)
 =\sum_{k\ge0}Y_k(x_1,\ldots,x_k)\frac{t^k}{k!}.
\end{equation}
In particular $Y_0=1$, $Y_1=x_1$,
$Y_2=x_1^2+x_2$, and $Y_3=x_1^3+3x_1x_2+x_3$.
Write
\begin{equation}\label{gs:kappa}
 \kappa_j(a,b)=
 \psi^{(j-1)}(1)-\psi^{(j-1)}(a)-\psi^{(j-1)}(b),\qquad j\ge1.
\end{equation}

\begin{theorem}[All Stieltjes indices]\label{gs:spectral-jets}
For every integer $k\ge0$,
\begin{align}
 &\sum_{n=0}^{\infty}\left[
 W_n(a,b;0)\,
 Y_k\bigl(-\psi(n+a+b),\ldots,-\psi^{(k-1)}(n+a+b)\bigr)
 -\frac{(-\log(n+c))^k}{n+c}\right]\notag\\
 &\hspace{18mm}=
 \frac{Y_{k+1}(\kappa_1,\ldots,\kappa_{k+1})}{k+1}
 -(-1)^k\gamma_k(c).\label{gs:all-index}
\end{align}
For $k=0$, the Bell factor in the summand is $1$ and has no arguments.
The series is absolutely convergent, locally uniformly in the shift
parameters.  Its summand is $O((1+\log n)^k/n^2)$ on real compact
parameter sets.
\end{theorem}

\begin{proof}
The $k$th $s$ derivative of $W_n$ at zero is the first term in brackets,
because the $j$th logarithmic derivative of its sole $s$-dependent gamma
factor is $-\psi^{(j-1)}(n+a+b)$.  On the other hand,
\[
 R_{a,b}^{(k+1)}(0)=Y_{k+1}(\kappa_1,\ldots,\kappa_{k+1}),
 \qquad
 \left.\frac{d^k}{ds^k}\frac{R_{a,b}(s)-1}{s}\right|_{s=0}
 =\frac{R_{a,b}^{(k+1)}(0)}{k+1}.
\]
Differentiate \eqref{gs:master-equation} using
\eqref{gs:stieltjes-convention}.  Normal convergence follows from
Theorem~\ref{gs:master}.  The explicit logarithmic summand estimate follows
by differentiating the locally uniform gamma-ratio expansion, or by
using the corresponding polygamma expansions.
\end{proof}

The first three identities are
\begin{align}
 \sum_{n\ge0}\left(W_n-\frac1{n+c}\right)
 &=\kappa_1-\gamma_0(c),\label{gs:k0}\\
 \sum_{n\ge0}\left(W_n\psi(n+a+b)
                   -\frac{\log(n+c)}{n+c}\right)
 &=-\frac{\kappa_1^2+\kappa_2}{2}-\gamma_1(c),\label{gs:k1}\\
 \sum_{n\ge0}\left(W_n\bigl(\psi(n+a+b)^2-\psi'(n+a+b)\bigr)
                   -\frac{\log^2(n+c)}{n+c}\right)
 &=\frac{\kappa_1^3+3\kappa_1\kappa_2+\kappa_3}{3}
   -\gamma_2(c),\label{gs:k2}
\end{align}
where $W_n=W_n(a,b;0)$.

\subsection{A complete mixed-parameter coefficient rule}

The finite-dimensional Bell expression has a direct multivariate
version.  Put $q_n=n+a+b$ and introduce formal increments $u,v,s$.
Define
\begin{align}
 \mathcal K_n(u,v,s)
 &=\sum_{j\ge1}\frac1{j!}
 \left[\psi^{(j-1)}(n+a)u^j+\psi^{(j-1)}(n+b)v^j
       -\psi^{(j-1)}(q_n)(u+v+s)^j\right],\label{gs:Kn}\\
 \mathcal R(u,v,s)
 &=\frac{\Gamma(1+s)\Gamma(a+u)\Gamma(b+v)}
 {\Gamma(a+u+s)\Gamma(b+v+s)}.\label{gs:mixed-R}
\end{align}
Its logarithm is the explicit formal series
\begin{align}
 \log\mathcal R(u,v,s)
 =\sum_{j\ge1}\frac1{j!}\bigl[&\psi^{(j-1)}(1)s^j
 +\psi^{(j-1)}(a)\{u^j-(u+s)^j\}\notag\\
 &+\psi^{(j-1)}(b)\{v^j-(v+s)^j\}\bigr].\label{gs:mixed-logR}
\end{align}
For fixed indices, only finitely many terms of these exponentials are
needed.

\begin{theorem}[All mixed jets]\label{gs:mixed-jets}
For all integers $p,q,k\ge0$,
\begin{align}
 &\sum_{n\ge0}\left[
 W_n(a,b;0)[u^pv^qs^k]e^{\mathcal K_n(u,v,s)}
 -\delta_{p0}\delta_{q0}\frac{(-\log(n+c))^k}{k!(n+c)}\right]
 \notag\\
 &\qquad=[u^pv^qs^{k+1}]\mathcal R(u,v,s)
 -\delta_{p0}\delta_{q0}\frac{(-1)^k\gamma_k(c)}{k!}.
 \label{gs:mixed-coefficient}
\end{align}
All series are ordinary absolutely convergent series.  Multiplying by
$p!q!k!$ converts this coefficient formula into the mixed derivative
formula.  When $p+q>0$, no logarithmic subtraction and no Stieltjes
constant remain.
\end{theorem}

\begin{proof}
Taylor expansion of the three nonconstant gamma factors in
$W_n(a+u,b+v;s)$ gives
$W_n(a+u,b+v;s)=W_n(a,b;0)e^{\mathcal K_n(u,v,s)}$.
Apply Theorem~\ref{gs:master} at $(a+u,b+v,c)$ and compare the indicated
coefficient.  The identity $\mathcal R(u,v,0)=1$ explains the shift from
$s^k$ to $s^{k+1}$ without any division ambiguity.  Normal convergence
justifies the coefficient extraction under the sum.
\end{proof}

One useful consequence is a cancellation at every positive mixed
tangential order:
\begin{equation}\label{gs:tangential}
 \sum_{n\ge0}\partial_a^p\partial_b^q W_n(a,b;0)=0,
 \qquad p,q\ge1.
\end{equation}
Indeed, the $s^1$ coefficient of $\mathcal R$ is
$\psi(1)-\psi(a+u)-\psi(b+v)$, which has no mixed $u,v$ terms.
The pure derivatives satisfy
\begin{equation}\label{gs:pure-tangential}
 \sum_{n\ge0}\partial_a^p W_n(a,b;0)=-\psi^{(p)}(a),
 \qquad p\ge1,
\end{equation}
and the analogous identity holds for $b$.

For an explicit quadratic identity, set
$A_n=\psi(n+a)-\psi(q_n)$ and
$B_n=\psi(n+b)-\psi(q_n)$.  Then
\begin{equation}\label{gs:quadratic}
 \sum_{n\ge0}W_n(a,b;0)
 \left(A_nB_n-\psi'(q_n)\right)=0.
\end{equation}
The first spectral lift is also closed:
\begin{align}
 \sum_{n\ge0}W_n(a,b;0)\bigl[&-\psi(q_n)
             \{A_nB_n-\psi'(q_n)\}
 -\psi'(q_n)(A_n+B_n)-\psi''(q_n)\bigr]
 =\psi'(a)\psi'(b).\label{gs:quadratic-lift}
\end{align}
This follows either from \eqref{gs:mixed-coefficient} at $(p,q,k)=(1,1,1)$
or by differentiating \eqref{gs:k1} with its original negative spectral
sign retained.  At $a=b=1$, \eqref{gs:quadratic} is the classical
$\zeta(2,1)=\zeta(3)$ identity written as a polygamma series.  We therefore
do not present that specialization as a new evaluation.

\subsection{Ordinary and central-binomial Stieltjes series}

Let $H_n^{(r)}=\sum_{j=1}^n j^{-r}$, with $H_0^{(r)}=0$, and write
$H_n=H_n^{(1)}$.  At $a=b=c=1$ the weight is $1/(n+1)$ and
$(\kappa_1,\kappa_2,\kappa_3)=(\gamma,-\zeta(2),2\zeta(3))$.
Equations~\eqref{gs:k1}--\eqref{gs:k2} become
\begin{align}
 \sum_{m=1}^{\infty}\frac{H_m-\gamma-\log m}{m}
 &=\frac{\zeta(2)-\gamma^2}{2}-\gamma_1,\label{gs:ordinary1}\\
 \sum_{m=1}^{\infty}
 \frac{(H_m-\gamma)^2+H_m^{(2)}-\zeta(2)-\log^2m}{m}
 &=\frac{\gamma^3-3\gamma\zeta(2)+2\zeta(3)}3-\gamma_2.
 \label{gs:ordinary2}
\end{align}
The all-order form follows from \eqref{gs:all-index}; for $j\ge2$ its
local Bell arguments are
$-\psi^{(j-1)}(m+1)=(-1)^{j+1}(j-1)!
\{\zeta(j)-H_m^{(j)}\}$.

A less elementary weight is obtained at $a=b=\tfrac12$, $c=1$.
Put
\[
 d_n=\frac{\binom{2n}{n}^2}{16^n},\qquad
 L=\log2,\qquad K=\gamma+4L.
\]
The duplication formula gives $W_n=\pi d_n$, while
\begin{equation}\label{gs:central-kappa}
 \kappa_1=K,\qquad
 \kappa_j=(3-2^{j+1})(-1)^j(j-1)!\zeta(j)\quad(j\ge2).
\end{equation}
Thus
\begin{align}
 \sum_{n=0}^{\infty}\left(\pi d_n-\frac1{n+1}\right)
 &=4\log2,\label{gs:central0}\\
 \sum_{n=0}^{\infty}\left(
 \pi d_n(H_n-\gamma)-\frac{\log(n+1)}{n+1}\right)
 &=\frac52\zeta(2)-\frac{K^2}2-\gamma_1,\label{gs:central1}\\
 \sum_{n=0}^{\infty}\left(
 \pi d_n\{(H_n-\gamma)^2+H_n^{(2)}-\zeta(2)\}
 -\frac{\log^2(n+1)}{n+1}\right)
 &=\frac{K^3}{3}-5K\zeta(2)+\frac{26}{3}\zeta(3)-\gamma_2.
 \label{gs:central2}
\end{align}
For every $k$, the same squared central-binomial weight gives
$\gamma_k$ through \eqref{gs:all-index} with
\eqref{gs:central-kappa}; no empirical relation search is involved.

The tangential identities yield convergent harmonic sums without a
Stieltjes subtraction.  Since
$\psi(n+\tfrac12)-\psi(n+1)=2(H_{2n}-H_n-L)$, they imply
\begin{align}
 \sum_{n=0}^{\infty}d_n(H_{2n}-H_n-L)&=-\frac\pi4,
 \label{gs:central-linear}\\
 \sum_{n=0}^{\infty}d_n\left[
 4(H_{2n}-H_n-L)^2+H_n^{(2)}-\zeta(2)\right]&=0.
 \label{gs:central-quadratic}
\end{align}
The brackets must remain grouped: some of their separately summed
pieces diverge.  The summands displayed above are $O(n^{-2})$.

\subsection{Polylogarithmic boundary values}

For $|z|<1$, define
\begin{align}
 \mathcal F_{a,b,c}(s;z)
 &=\sum_{n\ge0}
 \left(W_n(a,b;s)-(n+c)^{-1-s}\right)z^n\notag\\
 &=\frac{\Gamma(a)\Gamma(b)}{\Gamma(a+b+s)}
 {}_2F_1(a,b;a+b+s;z)-\Phi(z,1+s,c).
 \label{gs:polylog-generator}
\end{align}
At lower hypergeometric parameters that are nonpositive integers, the
first expression specifies the regularized meaning of the normalized
hypergeometric factor.  Uniform absolute convergence on $|z|\le1$,
locally in $\Re s>-1$, proves the exact boundary identity
\begin{equation}\label{gs:boundary}
 \lim_{z\to1,\ |z|\le1}\partial_a^p\partial_b^q\partial_s^k
 \mathcal F_{a,b,c}(s;z)
 =\partial_a^p\partial_b^q\partial_s^k\mathcal F_{a,b,c}(s).
\end{equation}
The limit is not restricted to a radial approach.  In particular, for
$c=1$ the Lerch factor is
\begin{equation}\label{gs:Li-bridge}
 \Phi(z,1+s,1)=\frac{\operatorname{Li}_{1+s}(z)}z,
\end{equation}
so \eqref{gs:all-index} is also a family of exact boundary constants
for polylogarithm order derivatives after an explicit hypergeometric
subtraction.

For rational $c=r/q$, $1\le r\le q$, let $\omega=e^{2\pi i/q}$ and,
for example, take $0<z<1$.  The root-of-unity filter gives
\begin{equation}\label{gs:rational-Lerch}
 \Phi(z,1+s,r/q)=q^s z^{-r/q}
 \sum_{j=0}^{q-1}\omega^{-jr}
 \operatorname{Li}_{1+s}(\omega^jz^{1/q}).
\end{equation}
Indeed the sum selects positive exponents congruent to $r$ modulo $q$;
after writing those exponents as $qn+r$, the prefactor becomes exactly
$(n+r/q)^{-1-s}$.  Finite translation reduces every positive rational
$c$ to this range.  All $s$ derivatives bring only finite binomial
combinations with powers of $\log q$, in addition to polylogarithm order
derivatives.

At the central-binomial point $s=0$, the first term in
\eqref{gs:polylog-generator} is a classical complete elliptic integral
$\pi\,{}_2F_1(\tfrac12,\tfrac12;1;z)$.
It is twice the complete elliptic integral with parameter $z$.
Consequently \eqref{gs:central0} is its logarithmic boundary constant,
whereas \eqref{gs:central1} and \eqref{gs:central2} describe the next
spectral layers of the same cancellation.

% Additional research fragment; requires the definitions in gauss_stieltjes_notes.tex.

\section{Successive subtractions and every negative balance}
\label{gs:higher-balance}

The first subtraction in Theorem~\ref{gs:master} has an exact extension
through every negative integer balance.  Its coefficients are finite
polynomials, and residue cancellation gives an identity among those
polynomials that is independent of the comparison shift.

Let $B_j(x)$ be the Bernoulli polynomials, defined by
\[
 \frac{te^{xt}}{e^t-1}=\sum_{j\ge0}B_j(x)\frac{t^j}{j!}.
\]
For $j\ge1$, set
\begin{align}
 D_j(a,b,s;c)=\frac{(-1)^{j+1}}{j(j+1)}
 \bigl[&B_{j+1}(a-c)+B_{j+1}(b-c)\notag\\
       &-B_{j+1}(a+b+s-c)-B_{j+1}(1-c)\bigr].
 \label{gs:Dj}
\end{align}
Define polynomials $A_j=A_j(a,b,s;c)$ by the formal identity
\begin{equation}\label{gs:Aj-definition}
 \sum_{j\ge0}A_jt^j
 =\exp\left(\sum_{j\ge1}D_jt^j\right).
\end{equation}
Equivalently,
\begin{equation}\label{gs:Aj-recursion}
 A_0=1,\qquad
 A_m=\frac1m\sum_{j=1}^m jD_jA_{m-j}\quad(m\ge1).
\end{equation}
This definition is purely finite for any fixed $m$.  The series in
\eqref{gs:Aj-definition} is used formally; it need not converge as an
infinite series in $t$.

\begin{theorem}[Exact continuation by successive coefficient subtraction]
\label{gs:successive-subtraction}
Fix an integer $M\ge0$.  On the domain
\[
 \Re a>0,\quad \Re b>0,\quad \Re c>0,\quad \Re s>-M-1,
\]
the series
\begin{equation}\label{gs:FM}
 \mathcal F^{[M]}_{a,b,c}(s)=
 \sum_{n=0}^{\infty}\left[
 W_n(a,b;s)-\sum_{j=0}^M
 A_j(a,b,s;c)(n+c)^{-1-s-j}\right]
\end{equation}
converges normally on compact subsets and is jointly holomorphic.  It
equals the holomorphic continuation of
\begin{equation}\label{gs:FM-closed}
 \frac{\Gamma(a)\Gamma(b)\Gamma(s)}
 {\Gamma(a+s)\Gamma(b+s)}
 -\sum_{j=0}^M A_j(a,b,s;c)\zeta(1+s+j,c).
\end{equation}
Every apparent pole at $s=0,-1,\ldots,-M$ is removable in this complete
expression.  For every integer $r\ge0$, the finite polynomial identity
\begin{equation}\label{gs:residue-rigidity}
 A_r(a,b,-r;c)
 =\frac{(-1)^r\Gamma(a)\Gamma(b)}
        {r!\,\Gamma(a-r)\Gamma(b-r)}
 =\frac{(-1)^r}{r!}(1-a)_r(1-b)_r
\end{equation}
holds.  In particular its left side is independent of $c$.
\end{theorem}

\begin{proof}
Write $x=n+c$.  The locally uniform logarithmic gamma expansion in a
fixed sector around the positive real axis gives
\begin{equation}\label{gs:log-asymptotic}
 \log W_n(a,b;s)
 =-(1+s)\log x+\sum_{j=1}^M\frac{D_j(a,b,s;c)}{x^j}
   +O(x^{-M-1}).
\end{equation}
For all sufficiently large $n$ the gamma factors have no zeros or poles
on the compact parameter set, so a logarithm may be chosen consistently
there.  The coefficient formula follows directly by substituting the
four shifts $a-c$, $b-c$, $a+b+s-c$, $1-c$ into the logarithmic gamma
expansion in Bernoulli polynomials.  Its leading powers combine to
$-(1+s)\log x$, while the constant terms cancel.
The remainder is uniform on compact parameter sets: all four shifts
are bounded, and all four large gamma arguments lie in a common closed
sector that avoids the negative real axis.

Exponentiating \eqref{gs:log-asymptotic} to finite order and using
\eqref{gs:Aj-definition} yields
\[
 W_n(a,b;s)=x^{-1-s}
 \left(\sum_{j=0}^M A_j(a,b,s;c)x^{-j}+O(x^{-M-1})\right).
\]
Thus the bracket in \eqref{gs:FM} is
$O(n^{-M-2-\Re s})$.  On a compact subset of the stated domain it is
$O(n^{-1-\eta})$ for some $\eta>0$.  This proves normal convergence and
joint holomorphy.  The finitely many small $n$ cause no difficulty:
$W_n$ is defined using the entire reciprocal gamma factor in
\eqref{gs:weight}, including at its zeros.

For $\Re s>0$, every series can be summed separately and Gauss
summation gives \eqref{gs:FM-closed}.  Its meromorphic continuation in
the larger half-plane has possible poles only at $s=0,-1,\ldots,-M$.
The identity theorem, first away from those points and then across
them, identifies it with the holomorphic left side.  All its residues
there must vanish.  At $s=-r$, only the term indexed by $j=r$ in the
zeta sum can have a pole, and its residue is $A_r(a,b,-r;c)$.
The gamma quotient has residue
\[
 \frac{(-1)^r\Gamma(a)\Gamma(b)}
 {r!\,\Gamma(a-r)\Gamma(b-r)}.
\]
This expression is interpreted through reciprocal gamma if either
denominator argument is a nonpositive integer; its value is then zero,
and the corresponding gamma pole is canceled.  Equating the two
residues proves the first equality in \eqref{gs:residue-rigidity}.
The gamma recurrence proves the polynomial expression in the second
equality.  Since $M$ can be any integer at least $r$, the result holds
for every $r$.
\end{proof}

The quotient in \eqref{gs:residue-rigidity} is holomorphic on the stated
$a,b$ domain even when $a-r$ or $b-r$ is a nonpositive integer.  The
Pochhammer expression makes this particularly transparent.  One must
not replace such a reciprocal gamma factor by an undefined division
by an infinite gamma value, nor differentiate a logarithm through one
of its zeros.

\subsection{An exact jet formula at a negative integer}

Fix $0\le r\le M$.  Suppress $a,b,c$ from the notation for $A_j$ and
write $A_j^{(k)}$ for its $k$th derivative in $s$.  Define the analytic
germ
\begin{equation}\label{gs:Rr}
 \mathcal R_r(\varepsilon)=
 \frac{\Gamma(1+\varepsilon)\Gamma(a)\Gamma(b)}
 {\displaystyle\prod_{h=1}^{r}(\varepsilon-h)}
 \frac1{\Gamma(a-r+\varepsilon)\Gamma(b-r+\varepsilon)}.
\end{equation}
The empty product is $1$.  This germ is analytic for sufficiently small
$\varepsilon$, also when either reciprocal gamma factor vanishes at
zero.  It equals $\varepsilon$ times the meromorphic gamma quotient
in \eqref{gs:FM-closed} evaluated at $s=-r+\varepsilon$.

\begin{corollary}[Every resonant spectral jet]\label{gs:resonant-jet}
For each integer $k\ge0$,
\begin{align}
 \left(\mathcal F^{[M]}_{a,b,c}\right)^{(k)}(-r)
 ={}&\frac{\mathcal R_r^{(k+1)}(0)-A_r^{(k+1)}(-r)}{k+1}
 \notag\\
 &-\sum_{\ell=0}^k\binom{k}{\ell}
 A_r^{(k-\ell)}(-r)(-1)^\ell\gamma_\ell(c)
 \notag\\
 &-\sum_{\substack{0\le j\le M\\j\ne r}}
 \sum_{\ell=0}^k\binom{k}{\ell}
 A_j^{(k-\ell)}(-r)\zeta^{(\ell)}(1-r+j,c).
 \label{gs:resonant-jet-formula}
\end{align}
The derivatives on the left are given by termwise derivatives of the
convergent brackets in \eqref{gs:FM}.  This formula remains valid at
all admissible $a,b,c$, including the parameter values at which a gamma
factor in an unregularized logarithmic expression would have a pole.
\end{corollary}

\begin{proof}
Put $Z(\varepsilon,c)=\zeta(1+\varepsilon,c)-1/\varepsilon$.
Equations~\eqref{gs:FM-closed} and \eqref{gs:Rr} give the exact germ
\begin{align*}
 \mathcal F^{[M]}(-r+\varepsilon)
 ={}&\frac{\mathcal R_r(\varepsilon)-A_r(-r+\varepsilon)}{\varepsilon}
 -A_r(-r+\varepsilon)Z(\varepsilon,c)\\
 &-\sum_{j\ne r}A_j(-r+\varepsilon)
 \zeta(1-r+j+\varepsilon,c).
\end{align*}
The numerator of the first term vanishes at zero by
\eqref{gs:residue-rigidity}.  Taking $k$ derivatives and using
\eqref{gs:stieltjes-convention} gives the formula.  Normal convergence
on slightly larger compact polydiscs justifies its series version.
\end{proof}

If neither $a-r$ nor $b-r$ is a nonpositive integer, put
$Q_r=A_r(a,b,-r;c)$.  The finite value $k=0$ can then be written
\begin{align}
 \mathcal F^{[M]}(-r)
 ={}&Q_r\bigl(\psi(r+1)-\psi(a-r)-\psi(b-r)\bigr)
 -A_r'(-r)-Q_r\gamma_0(c)\notag\\
 &-\sum_{j\ne r}A_j(-r)\zeta(1-r+j,c).
 \label{gs:resonant-finite}
\end{align}
At the excluded logarithmic parameter values the analytic germ
\eqref{gs:Rr}, rather than an informal $0\cdot\infty$ convention,
specifies the value.  Higher derivatives of this germ can also be
written as Bell polynomials.  When $Q_r\ne0$, its logarithmic arguments
are
\[
 \psi^{(j-1)}(1)+(j-1)!H_r^{(j)}
 -\psi^{(j-1)}(a-r)-\psi^{(j-1)}(b-r),\qquad j\ge1.
\]

\subsection{An exact central-binomial identity at balance minus one}

Take $a=b=1/2$, $c=1$, $M=r=1$.  Formula~\eqref{gs:Dj} gives
\[
 A_1(s)=\frac34+\frac s2-\frac{s^2}{2},\qquad
 A_1(-1)=-\frac14,\qquad A_1'(-1)=\frac32.
\]
Moreover,
\[
 W_n(1/2,1/2;-1)=\pi n\frac{\binom{2n}{n}^2}{16^n}.
\]
At $n=0$ this equality has its literal value $0$: the original gamma
ratio contains $1/\Gamma(0)=0$, interpreted through the entire
reciprocal gamma function.  The rest of the $n=0$ bracket remains
$-1+1/4=-3/4$ and is included in the sum.

Using $\psi(2)=1-\gamma$,
$\psi(-1/2)=2-\gamma-2\log2$,
$\zeta(0)=-1/2$, and $\gamma_0(1)=\gamma$ in
\eqref{gs:resonant-finite} yields the exact identity
\begin{equation}\label{gs:central-minus-one}
 \boxed{\displaystyle
 \sum_{n=0}^{\infty}\left[
 \pi n\frac{\binom{2n}{n}^2}{16^n}
 -1+\frac1{4(n+1)}\right]
 =-\log2-\frac14.}
\end{equation}
The bracket is $O(n^{-2})$.  Its three displayed pieces must be summed
together; the individual series do not converge.

The numerical verification evaluates this bracket directly, then uses
the independent gamma-ratio asymptotic tail
\[
 \sum_{j=2}^K A_j(1/2,1/2,-1;1)\zeta(j,N+1).
\]
It checks stability under two choices of $N,K$.  The same script verifies
\eqref{gs:residue-rigidity} symbolically for $r=0,1,2,3,4$ with all of
$a,b,c$ indeterminate.  These computational checks support, and are
separate from, the all-order proof above.

\section{Correcting and proving a historical family of harmonic conjectures}
\label{gs:choi-section}

The paper of Choi \cite{gs:Choi} uses the standard definition
$H_n=\sum_{j=1}^n1/j$.  Its printed page~6 contains Conjecture~4,
Corollary~5, and Conjecture~6.  The formulas can be checked and extended
uniformly by the following generating identity.  This discussion
concerns those displayed statements, and does not assert that they
remain open or that no correction has appeared elsewhere.

\begin{theorem}[A two-parameter harmonic generating identity]
\label{gs:choi-master}
For $u,v$ in a neighborhood of zero,
\begin{equation}\label{gs:choi-G}
 \sum_{n=0}^{\infty}
 \frac{(1+u)_n(1+v)_n}{(n!)^2(n+1)(n+2)}
 =\frac{\Gamma(1-u-v)}{\Gamma(2-u)\Gamma(2-v)}.
\end{equation}
Define
\begin{equation}\label{gs:En}
 E_r(n)=Y_r\bigl(H_n,-H_n^{(2)},2H_n^{(3)},\ldots,
              (-1)^{r-1}(r-1)!H_n^{(r)}\bigr),\qquad E_0(n)=1.
\end{equation}
Then
\begin{align}
 \sum_{n=1}^{\infty}\frac{E_r(n)}{(n+1)(n+2)}
 &=r!,&&r\ge1,\label{gs:choi-pure}\\
 \sum_{n=1}^{\infty}\frac{H_nE_r(n)}{(n+1)(n+2)}
 &=r!\left(1+\sum_{j=2}^{r+1}\zeta(j)\right),&&r\ge0.
 \label{gs:choi-mixed}
\end{align}
Every mixed sum with $E_r(n)E_t(n)$ is likewise given by the finite
Taylor coefficient $r!t![u^rv^t]$ of the gamma quotient in
\eqref{gs:choi-G}, except that the $r=t=0$ formula must retain the
$n=0$ term or subtract its value $1/2$.
\end{theorem}

\begin{proof}
Since $(3)_n=n!(n+1)(n+2)/2$, the series in
\eqref{gs:choi-G} is
$\tfrac12{}_2F_1(1+u,1+v;3;1)$.
Gauss summation proves the stated gamma quotient.  Uniform convergence
and differentiation follow, for instance, from
$O(n^{-2+\Re(u+v)})$ on a small closed polydisc around zero.
The coefficient identity
\[
 \frac{(1+u)_n}{n!}
 =\prod_{j=1}^n\left(1+\frac uj\right)
 =\sum_{r\ge0}\frac{E_r(n)}{r!}u^r
\]
follows from \eqref{gs:Bell-definition}; it also shows that $E_r(n)$ is
$r!$ times an elementary symmetric polynomial in $1,1/2,\ldots,1/n$.

Set $v=0$ in \eqref{gs:choi-G}.  Its right side becomes $1/(1-u)$,
which proves \eqref{gs:choi-pure}.  For the second identity,
differentiate in $v$ at zero to obtain
\[
 \left.\partial_v
 \frac{\Gamma(1-u-v)}{\Gamma(2-u)\Gamma(2-v)}\right|_{v=0}
 =\frac{\psi(2)-\psi(1-u)}{1-u}
 =\frac{1+\sum_{j\ge1}\zeta(j+1)u^j}{1-u}.
\]
Taking the coefficient of $u^r$ proves \eqref{gs:choi-mixed}.
The zero-th summand has zero derivative in every case used in those
two identities, which explains their lower limit $n=1$.
\end{proof}

In the notation of Theorem~\ref{harm:thm:gamma-generating},
$E_r(n)=r!E_{n,r}$.  The pure identity also follows independently from
the resolvent difference at $a=1,b=2$, since $F(1,y)=1$ and
$F(2,y)=(1-y)^{-1}$.  The product with $H_n$ requires the two-parameter
coefficient in the theorem.

The theorem gives the intended fifth- and sixth-order cases of
Conjecture~4 through \eqref{gs:choi-pure}.  It also gives corrected
versions of all six displayed formulas (41)--(46) through
\eqref{gs:choi-mixed}.  The precise corrections are:
\begin{enumerate}
 \item The right sides in (41)--(46) must be divided by two; for (46),
 the numerator correction below is also required.  With those
 corrections, the multiplier in the row indexed by $r$ is $r!$, for
 $r=1,\ldots,6$, instead of the printed $2r!$.
 \item In the second identity of (40), the final harmonic term in the
 sixth-order Bell polynomial must be $-120H_n^{(6)}$; the printed
 version has $-120H_n^{(5)}$.
 \item In (46), the term belonging to the partition $2+2+1+1$ must
 be $45H_n^3(H_n^{(2)})^2$ after the outer factor $H_n$ is included.
 The printed formula has $45H_n^3(H_n)^2$ at that position.
\end{enumerate}
These readings were checked against the publisher's PDF, not inferred
from its mathematical-text extraction.  The useful correction is the
unambiguous Bell-polynomial statement, which also fixes all higher
orders without enumerating partitions individually.

The simplest factor-two error admits a short proof independent of
hypergeometric differentiation.  By summation by parts and Tonelli's
theorem for the nonnegative terms,
\begin{align*}
 \sum_{n=1}^{\infty}\frac{H_n^2}{(n+1)(n+2)}
 &=\sum_{k=1}^{\infty}\frac{H_k^2-H_{k-1}^2}{k+1}\\
 &=2\sum_{k=1}^{\infty}\frac{H_{k-1}}{k(k+1)}
   +\sum_{k=1}^{\infty}\frac1{k^2(k+1)}\\
 &=2\sum_{j=1}^{\infty}\frac1{j(j+1)}+\zeta(2)-1
 =1+\zeta(2).
\end{align*}
This disproves the printed value $2(1+\zeta(2))$ in (41) with its
stated definitions.  It also isolates the origin of the normalization:
the hypergeometric function is twice the left side of
\eqref{gs:choi-G}, rather than equal to it.

